\section{Scenario Analyses and Stakeholder Impact}

% tighter tables for this appendix
\setlength{\tabcolsep}{3pt}
\newcolumntype{Q}[1]{>{\raggedright\arraybackslash\hspace{0pt}\sloppy}p{#1}}
\setlength{\LTpre}{3pt}
\setlength{\LTpost}{3pt}

%----------------------------------------------------------------------
\subsection*{Flagship scenario – Terry stop}

\noindent\textbf{Context.} Evening patrol on a commercial strip. Officer pulls over a driver for a broken taillight, requests conversation beside the vehicle, and begins casual questioning.

\medskip\noindent\textbf{Micro‑transcript.}
\begin{quote}
\textbf{Civilian:} “I invoke my right to remain silent. I invoke my right to a lawyer. I don’t consent to any searches. I want to leave. Am I free to go?”\\[2pt]
\textbf{Officer:} “Understood. Please wait here while I verify your license and registration.”\\[2pt]
( License check clears; officer issues fix‑it ticket and releases driver. )
\end{quote}

\medskip\noindent\textbf{Outcomes.}

\begin{longtable}{@{}P{0.20\textwidth} Q{0.28\textwidth} Q{0.52\textwidth}@{}}
\toprule
\textbf{Actor} & \textbf{Immediate effect} & \textbf{Down‑stream benefit}\\
\midrule
\endfirsthead
\toprule
\textbf{Actor} & \textbf{Immediate effect} & \textbf{Down‑stream benefit}\\
\midrule
\endhead
\midrule
\multicolumn{3}{r}{\emph{Continued on next page}}\\
\midrule
\endfoot
\bottomrule
\endlastfoot

Officer &
Maintains procedural clarity; avoids unrecorded questioning. &
Lower liability exposure; no suppression risk if citation contested.\\[6pt]

Civilian &
Rights asserted without escalation. &
Clear evidence trail if dispute arises.\\[6pt]

Supervisor &
Body‑cam shows textbook compliance. &
Easier performance review; supports training benchmarks.\\[6pt]

Prosecutor &
(If case later filed) invocation time‑stamp defines questioning window. &
Reduced suppression‑motion workload.\\
\end{longtable}

\medskip\noindent\textbf{Failure‑mode variant.} Officer ignores the Invocation and continues questioning. Court later suppresses statements; Internal Affairs opens review for policy breach. Administrative cost ≈ new training session + IA hours; suppression cost = dismissed charge.

%----------------------------------------------------------------------
\subsection*{Stakeholder impact matrix}

\begin{longtable}{@{}P{0.20\textwidth} P{0.28\textwidth} Q{0.18\textwidth} Q{0.34\textwidth}@{}}
\toprule
\textbf{Stakeholder} & \textbf{Primary benefit} & \textbf{Trade‑off / cost} & \textbf{Training / resource need}\\
\midrule
\endfirsthead
\toprule
\textbf{Stakeholder} & \textbf{Primary benefit} & \textbf{Trade‑off / cost} & \textbf{Training / resource need}\\
\midrule
\endhead
\midrule
\multicolumn{4}{r}{\emph{Continued on next page}}\\
\midrule
\endfoot
\bottomrule
\endlastfoot

Civilians &
Clarity of rights; reduced stress during stops. &
None if trained; confusion if untrained. &
Community outreach; pocket cards.\\[6pt]

Front‑line officers &
Shield from perjury claims; clear stop boundaries. &
Less reliance on spontaneous admissions. &
Roll‑call briefing; updated policy memo.\\[6pt]

Supervisors &
Objective body‑cam benchmark for compliance. &
Must track invocation‑adherence metrics. &
Audit‑template updates.\\[6pt]

Prosecutors &
Cleaner evidentiary record; fewer suppression fights. &
Some cases require additional corroboration. &
Charging‑decision checklist tweak.\\[6pt]

Judges &
Narrower procedural disputes; faster hearings. &
Potential docket shift (fewer plea bargains). &
None.\\[6pt]

Elected officials &
Greater institutional legitimacy; reduced settlement payouts. &
Short‑term dip in conviction‑volume statistics. &
Public‑messaging toolkit.\\
\end{longtable}

%----------------------------------------------------------------------
\subsection*{Compatibility sidebar – stop‑and‑identify laws}

Line 4 (“I want to leave. Am I free to go?”) addresses detention status; it does not relieve statutory duties to provide identification where required. Courts treat ID requirements as distinct from speech‑based invocation (see \textit{Hiibel v.\ Sixth Judicial District Court} (2004)).

%----------------------------------------------------------------------
\subsection*{Transition‑cost box – items to budget once pilots begin}

\begin{itemize}\setlength\itemsep{2pt}
  \item Officer in‑service module (≈ 30 min)
  \item Wallet‑card printing / digital download host
  \item Policy‑memo revision and distribution
  \item Body‑cam audit‑rubric adjustment
\end{itemize}

Dollar values depend on jurisdictional scale and will be estimated after pilot design.

%----------------------------------------------------------------------
\subsection*{Evidence gap \& research agenda}

\begin{itemize}\setlength\itemsep{2pt}
  \item What proportion of body‑cam footage already contains partial or failed invocations?
  \item Does standardized invocation reduce the average number of suppression‑motion hearings per 1 000 arrests?
  \item What is the cost‑to‑savings ratio of officer training versus civil‑liability payouts over a three‑year horizon?
\end{itemize}

Researchers, agencies, or unions interested in collaborating on pilot studies are invited to make contact at \texttt{info@invocationofrights.org}.

\medskip
\noindent\textbf{Pilot concept (draft).} A two‑arm mock‑stop study will measure (a) recall accuracy of the four‑line script after a single 20‑minute training session, (b) officer perception of cooperation, and (c) suppression‑motion outcomes in simulated hearings. A 2‑page protocol is available at \url{https://invocationofrights.org/pilot}; prospective partners are invited to review and co‑design field trials.

\medskip
\noindent Appendix D reflects preliminary analysis; quantitative validation will follow field trials.
